\scanpage{335}\mspage{251}
% l. 1–3
\deleted{\unsure{et}} et un système de flèches
\[
 F(j)\unsure{\mathbin{-}|\longrightarrow} X(j)\qquad (j\in\operatorname{Ob}J).
\]
\deleted{\unsure{Alors}} \inserted{alors, $G$ et} le système de flèches (4.10.5.1) peut être choisi de telle façon que les flèches composées
% l. 4–8
\[
 F(j)\unsure{\mathbin{-}|\longrightarrow} X(j)\unsure{\mathbin{-}|\longrightarrow}G(j)
\]
définissent un hom de sections
\[
 F\longrightarrow G
\]
\unsure{(de la} catégorie $\underline{\Gamma}(\unsure{\mathcal E|\mathscr E}/J)$).

% slanted
\margin{\uline{NB} Les conditions $\mathrm F1$ et $\unsure{\mathrm F'1|\mathrm F1}$ équivalent à dire \deleted{que}\inserted{que} le foncteur envisagé satisfait les conditions $\mathrm F1$ et $\mathrm F3$, ou \unsure{encore|mieux} \deleted{que} (en vertu de \unsure{4.9.6.3}, $1^\circ$) que $\underline{\Gamma}(\deleted{\unsure{\mathscr D|\mathcal E}}\mathcal E/J)$ est \uline{filtrante} (sans plus), et le foncteur \inserted{fidèle} (4.10.3.1) cofinal : Cela implique déjà que $\underline{\Gamma}(\mathcal E/J)$ est grde devant $\unsure{\pi|\kappa}$\unsure{.}}

% l. 9–12
\uline{Dém} Il faut voir que $\mathrm F1\mathbin{\&}\mathrm F'1$ équivaut à $\mathrm F1\mathbin{\&}\mathrm F3_{\pi}$. L'implication $\Longleftarrow$ est \inserted{$\pm$} immédiate, c'est de $\Longrightarrow$ que nous avons besoin, donc
\[
 \mathrm F1\mathbin{\&}\mathrm F'1\Longrightarrow\mathrm F3_{\pi}.
\]

% l. 13–19
Pour ceci, \deleted{rappelons} réénonçons $\mathrm F3_{\pi}$ dans le contexte présent. On se donne une famille $(F_{\alpha})_{\alpha\in A}$, $\operatorname{card}A\leq\pi$, d'objets de $\underline{\Gamma}(\mathcal E/J)$, un objet $F$ de la $\overline{\mathrm m}$ \unsure{catégor.|catégorie}, et \deleted{pour} tout $j$ dans $\operatorname{Ob}J$, un système de flèches
\[
 F_{\alpha}(j)\xrightarrow{u_{\alpha}(j)} F(j),
\]
mais sans supposer que pour $\alpha$ fixé, $j$ variable,
% archive sheet number: 334
